\chapter{The mathematical toolkit: vectors, projection, covariance, eigenvectors, PCA and whitening}\label{ch:linalg}

\begin{onething}
A pixel's spectrum is an arrow in a space with one axis per band, and almost every calculation
in this book is one of a handful of operations on such arrows: measuring length and angle,
projecting one arrow onto another, and turning a cloud of arrows so that it is round instead of
stretched. That last operation, whitening, is what the statistical detectors of
\cref{ch:rx,ch:localrx} and the residual score of \cref{ch:residualmaha} are built on.
\end{onething}

\section{Why a chapter of linear algebra}\label{sec:linalg:why}

The next two chapters build the classical anomaly detectors: the RX detector, its local version,
and MNF-RX. The chapter after them in Part VI builds the score that Indradhanu ships with,
residual Mahalanobis. All four do the same three things to a cloud of spectra:
\begin{enumerate}
  \item find the centre of the cloud (a \emph{mean});
  \item find the directions in which the cloud is stretched, and how much (a \emph{covariance}
    and its \emph{eigenvectors});
  \item rescale every direction so that the cloud becomes round (\emph{whitening}), and measure
    how far a pixel sits from the centre of the round cloud.
\end{enumerate}

You could try to skip the machinery and score each band on its own: subtract the band's average,
divide by its spread, add up. \Cref{ch:rx} shows in detail why that fails, but the short version
is this. Bands are not independent. Two neighbouring bands of a vegetation spectrum rise and fall
together across a scene, so a pixel that is a little high in both is ordinary, and a pixel that
is high in one and low in the other is strange, even if neither value alone is extreme. Seeing
the second case needs the arithmetic of how bands vary \emph{together}, which is covariance, and
the geometry of directions in band space, which is this chapter.

The anomaly-detection course gives this material as a warm-up of five operations (dot product,
norm, angle, matrix product, softmax). This chapter keeps the first four, raises them to a
derivation, and adds what the detectors need: projection, transpose, the inverse, covariance,
eigenvectors, PCA and whitening. Softmax and the 1 $\times$ 1 convolution belong to the neural
network chapters; they are taught in \cref{ch:shapes,ch:attention}.

\begin{analogy}
Think of exam marks. A raw mark of 70 means little until you know the class average and how
spread out the class was: 70 in a test where everyone scored between 68 and 72 is unremarkable,
and 70 in a test where the average was 40 with a spread of 10 is three spreads above the
average. Converting marks to ``spreads above average'' is the one-dimensional version of
whitening. With two subjects that are correlated (students good at algebra tend to be good at
geometry) you cannot convert each subject separately; you must first turn the axes to the
directions in which the class actually varies, then convert. That turning-then-converting is
what this chapter builds.
\end{analogy}

\section{Vectors, the dot product, length and angle}\label{sec:linalg:vectors}

\begin{definition}[Vector]\index{vector}
A \emph{vector} with $C$ entries is an ordered list of $C$ numbers, written as a column
\[
\vect{x} = \begin{pmatrix} x_1 \\ x_2 \\ \vdots \\ x_C \end{pmatrix}.
\]
The set of all such lists is written $\R^C$. You can picture a vector as an arrow from the
origin to the point whose coordinates are its entries.
\end{definition}

A pixel's spectrum is a vector: $x_1$ is its reflectance in band 1, $x_2$ in band 2, and so on.
On the 165-band common grid of \cref{ch:sensors} a spectrum lives in $\R^{165}$. You cannot draw
165 axes, but every formula below works the same in 2, 3 or 165 dimensions, so we draw the 2-D
case and trust the algebra for the rest.

Two operations are defined entry by entry. Adding two vectors adds their entries,
$(\vect{x}+\vect{y})_i = x_i + y_i$; multiplying a vector by a number $t$ (a \emph{scalar})
multiplies every entry, $(t\vect{x})_i = t x_i$. Geometrically, $t\vect{x}$ is the same arrow
stretched by $t$; for a spectrum it is the same shape at $t$ times the brightness.

\subsection{The dot product}

\begin{definition}[Dot product]\index{dot product}
The \emph{dot product} (or inner product) of two vectors of the same length is the number
\begin{equation}\label{eq:linalg:dot}
\ip{\vect{a}}{\vect{b}} = \vect{a}\T\vect{b} = a_1 b_1 + a_2 b_2 + \cdots + a_C b_C
= \sum_{i=1}^{C} a_i b_i .
\end{equation}
\end{definition}

The notation $\vect{a}\T\vect{b}$ anticipates \cref{sec:linalg:matrices}: a row times a column.
The result is a single number, not a vector. Three properties follow directly from the
definition, and you will use all three:
\begin{itemize}
  \item symmetry: $\vect{a}\T\vect{b} = \vect{b}\T\vect{a}$, because $a_i b_i = b_i a_i$;
  \item scaling: $(t\vect{a})\T\vect{b} = t(\vect{a}\T\vect{b})$, because $t$ factors out of
    every term of the sum;
  \item adding: $(\vect{a}+\vect{c})\T\vect{b} = \vect{a}\T\vect{b} + \vect{c}\T\vect{b}$,
    because $(a_i + c_i)b_i = a_i b_i + c_i b_i$.
\end{itemize}

\begin{worked}[label={we:linalg:dot}]{The dot product of two 3-band spectra}
Take the two spectra of the course's warm-up, $\vect{a} = (0.1, 0.2, 0.3)$ and
$\vect{b} = (0.3, 0.2, 0.1)$: one rises towards the long wavelengths, the other falls.
Do this on paper. Multiply position by position:
\begin{align*}
\text{band 1:}\quad 0.1 \times 0.3 &= 0.03\\
\text{band 2:}\quad 0.2 \times 0.2 &= 0.04\\
\text{band 3:}\quad 0.3 \times 0.1 &= 0.03
\end{align*}
then add:
\[
\vect{a}\T\vect{b} = 0.03 + 0.04 + 0.03 = 0.10 .
\]
\end{worked}

\subsection{Length}

\begin{definition}[Norm]\index{norm}
The \emph{length} or \emph{Euclidean norm} of a vector is
\begin{equation}\label{eq:linalg:norm}
\norm{\vect{a}} = \sqrt{\vect{a}\T\vect{a}} = \sqrt{a_1^2 + a_2^2 + \cdots + a_C^2}.
\end{equation}
\end{definition}

In two dimensions this is Pythagoras: the arrow to $(3, 4)$ is the hypotenuse of a right
triangle with sides 3 and 4, so its length is $\sqrt{9 + 16} = 5$. In more dimensions you apply
Pythagoras repeatedly, one axis at a time, and get the same formula. For a spectrum, the length
is essentially its overall brightness.

\begin{worked}[label={we:linalg:norm}]{The length of a spectrum}
For $\vect{a} = (0.1, 0.2, 0.3)$, square each entry:
\begin{align*}
0.1^2 &= 0.01\\
0.2^2 &= 0.04\\
0.3^2 &= 0.09
\end{align*}
add them, $0.01 + 0.04 + 0.09 = 0.14$, and take the square root:
$\norm{\vect{a}} = \sqrt{0.14} = 0.37417$.
The spectrum $\vect{b} = (0.3, 0.2, 0.1)$ has the same three squares in a different order, so
$\norm{\vect{b}} = 0.37417$ as well: the order of the bands does not affect the length.
\end{worked}

\subsection{Angle}

The dot product and the lengths together give the angle between two arrows.

\begin{definition}[Angle between vectors]\index{angle between vectors}
For non-zero $\vect{a}$ and $\vect{b}$ the angle $\theta$ between them satisfies
\begin{equation}\label{eq:linalg:cos}
\cos\theta = \frac{\vect{a}\T\vect{b}}{\norm{\vect{a}}\,\norm{\vect{b}}},
\qquad 0 \le \theta \le \pi .
\end{equation}
\end{definition}

Where does \cref{eq:linalg:cos} come from? Draw the triangle with sides $\vect{a}$, $\vect{b}$
and $\vect{a}-\vect{b}$. The school \emph{law of cosines} says the third side satisfies
\[
\norm{\vect{a}-\vect{b}}^2 = \norm{\vect{a}}^2 + \norm{\vect{b}}^2
  - 2\norm{\vect{a}}\norm{\vect{b}}\cos\theta .
\]
Now compute the left side with the dot product instead. Expand the square using the adding and
symmetry properties:
\begin{align*}
\norm{\vect{a}-\vect{b}}^2
&= (\vect{a}-\vect{b})\T(\vect{a}-\vect{b})
  && \text{definition of the norm}\\
&= \vect{a}\T\vect{a} - \vect{a}\T\vect{b} - \vect{b}\T\vect{a} + \vect{b}\T\vect{b}
  && \text{multiply out, term by term}\\
&= \norm{\vect{a}}^2 - 2\,\vect{a}\T\vect{b} + \norm{\vect{b}}^2
  && \text{symmetry: } \vect{b}\T\vect{a} = \vect{a}\T\vect{b}.
\end{align*}
Set the two expressions for $\norm{\vect{a}-\vect{b}}^2$ equal. The $\norm{\vect{a}}^2$ and
$\norm{\vect{b}}^2$ terms cancel, leaving
$-2\,\vect{a}\T\vect{b} = -2\norm{\vect{a}}\norm{\vect{b}}\cos\theta$. Divide both sides by
$-2\norm{\vect{a}}\norm{\vect{b}}$ and you have \cref{eq:linalg:cos}.

\begin{worked}[label={we:linalg:angle}]{The angle between two spectra}
Every piece is already computed in \cref{we:linalg:dot,we:linalg:norm}:
\begin{align*}
\norm{\vect{a}}\,\norm{\vect{b}} &= 0.37417 \times 0.37417 = 0.14\\
\cos\theta &= \frac{0.10}{0.14} = 0.71429\\
\theta &= \arccos(0.71429) = 0.7752\ \text{rad}.
\end{align*}
To convert radians to degrees multiply by $180/\pi$:
$0.7752 \times 180 / 3.14159 = 44.4^\circ$. The two spectra point in noticeably different directions.

Now replace $\vect{b}$ by $2\vect{a} = (0.2, 0.4, 0.6)$, the same shape at twice the brightness:
\begin{align*}
\vect{a}\T(2\vect{a}) &= 0.1\times0.2 + 0.2\times0.4 + 0.3\times0.6 = 0.02 + 0.08 + 0.18 = 0.28\\
\norm{2\vect{a}} &= \sqrt{0.04 + 0.16 + 0.36} = \sqrt{0.56} = 0.74833\\
\cos\theta &= \frac{0.28}{0.37417 \times 0.74833} = \frac{0.28}{0.28} = 1\\
\theta &= \arccos(1) = 0 .
\end{align*}
The angle is exactly zero although every entry differs. Angle measures shape and is blind to
brightness. That is the spectral angle of \cref{ch:losses}, and the reason the training loss
pairs it with an L1 term that does notice brightness.
\end{worked}

Two arrows with $\vect{a}\T\vect{b} = 0$ have $\cos\theta = 0$, so $\theta = \pi/2$: they are at
right angles. The word for this is \emph{orthogonal}\index{orthogonal}. A vector of length 1 is
a \emph{unit vector}\index{unit vector}; dividing any non-zero vector by its length makes one,
$\vect{u} = \vect{a}/\norm{\vect{a}}$.

\section{Projection: the best multiple of one vector}\label{sec:linalg:projection}

The single most useful operation in this book is \emph{projection}: given a vector $\vect{x}$
and a direction $\vect{u}$, find the multiple $t\vect{u}$ that comes closest to $\vect{x}$. You
will meet it again as ``the coordinate of a pixel along a principal direction'' (this chapter),
``the part of a residual inside the scene's main directions'' (the foreign fraction of
\cref{ch:benchmarks}), and as the heart of PCA.

\subsection{Derivation by minimising a squared distance}

Write the squared distance from $\vect{x}$ to $t\vect{u}$ as a function of the unknown number
$t$:
\[
f(t) = \norm{\vect{x} - t\vect{u}}^2 .
\]
Expand it, exactly as in the law-of-cosines derivation above:
\begin{align}
f(t) &= (\vect{x} - t\vect{u})\T(\vect{x} - t\vect{u})
  && \text{definition of the norm} \notag\\
&= \vect{x}\T\vect{x} - t\,\vect{x}\T\vect{u} - t\,\vect{u}\T\vect{x} + t^2\,\vect{u}\T\vect{u}
  && \text{multiply out; $t$ factors out of each term} \notag\\
&= \vect{x}\T\vect{x} - 2t\,\vect{u}\T\vect{x} + t^2\,\vect{u}\T\vect{u}
  && \text{symmetry of the dot product.} \label{eq:linalg:ft}
\end{align}
This is a parabola in $t$: a constant, a term linear in $t$, and $t^2$ times the positive number
$\vect{u}\T\vect{u}$. A parabola that opens upwards has its lowest point where its slope is
zero. Differentiate with respect to $t$ (the constant $\vect{x}\T\vect{x}$ has slope 0; the
derivative of $-2t\,c$ is $-2c$; the derivative of $t^2 c$ is $2tc$):
\[
\frac{df}{dt} = -2\,\vect{u}\T\vect{x} + 2t\,\vect{u}\T\vect{u}.
\]
Set the slope to zero and solve for $t$:
\begin{align}
-2\,\vect{u}\T\vect{x} + 2t^\star\,\vect{u}\T\vect{u} &= 0 \notag\\
t^\star\,\vect{u}\T\vect{u} &= \vect{u}\T\vect{x}
  && \text{add $2\vect{u}\T\vect{x}$ to both sides, divide by 2} \notag\\
t^\star &= \frac{\vect{u}\T\vect{x}}{\vect{u}\T\vect{u}}
  && \text{divide by $\vect{u}\T\vect{u} > 0$.} \label{eq:linalg:tstar}
\end{align}

\begin{definition}[Projection onto a vector]\index{projection}
The \emph{projection} of $\vect{x}$ onto the direction of $\vect{u}$ is
\begin{equation}\label{eq:linalg:proj}
\vect{p} = t^\star\vect{u} = \frac{\vect{u}\T\vect{x}}{\vect{u}\T\vect{u}}\,\vect{u},
\end{equation}
and the \emph{residual} of the projection is $\vect{r} = \vect{x} - \vect{p}$. If $\vect{u}$ is a
unit vector, $\vect{u}\T\vect{u} = 1$ and the formula shortens to
$\vect{p} = (\vect{u}\T\vect{x})\,\vect{u}$: the number $\vect{u}\T\vect{x}$ is then the
\emph{coordinate} of $\vect{x}$ along $\vect{u}$.
\end{definition}

The residual is always orthogonal to $\vect{u}$. Check it with the properties of the dot product:
\[
\vect{u}\T\vect{r} = \vect{u}\T\vect{x} - t^\star\,\vect{u}\T\vect{u}
= \vect{u}\T\vect{x} - \frac{\vect{u}\T\vect{x}}{\vect{u}\T\vect{u}}\,\vect{u}\T\vect{u}
= \vect{u}\T\vect{x} - \vect{u}\T\vect{x} = 0 .
\]
So $\vect{x} = \vect{p} + \vect{r}$ splits $\vect{x}$ into a part along $\vect{u}$ and a part at
right angles to it, and by Pythagoras $\norm{\vect{x}}^2 = \norm{\vect{p}}^2 + \norm{\vect{r}}^2$.

\begin{figure}[htbp]
\centering
\begin{tikzpicture}[scale=1.3,>=Stealth]
  \draw[->,gray] (-0.3,0) -- (3.8,0) node[right] {band 1};
  \draw[->,gray] (0,-0.3) -- (0,3.0) node[above] {band 2};
  \draw[spat,dashed] (-0.3,-0.3) -- (2.9,2.9) node[above right] {line $t\vect{u}$};
  \draw[->,very thick,spat] (0,0) -- (1,1) node[midway,above left] {$\vect{u}$};
  \draw[->,very thick,spec] (0,0) -- (3,1) node[right] {$\vect{x} = (3,1)$};
  \draw[->,very thick,attn] (0,0) -- (2,2) node[above left] {$\vect{p} = (2,2)$};
  \draw[->,thick,loss] (2,2) -- (3,1) node[midway,above right] {$\vect{r} = (1,-1)$};
  \draw[loss] ($(2,2)+(0.15,-0.15)$) -- ++(0.15,0.15) -- ++(-0.15,0.15);
  \fill[maskgrey] (0.95,0.95) circle (1.2pt);
  \fill[maskgrey] (2.85,2.85) circle (1.2pt);
  \draw[maskgrey,thin] (3,1) -- (0.95,0.95);
  \draw[maskgrey,thin] (3,1) -- (2.85,2.85);
\end{tikzpicture}
\caption{Projecting $\vect{x} = (3,1)$ onto the direction $\vect{u} = (1,1)$. The closest point of
the line $t\vect{u}$ is $\vect{p} = 2\vect{u} = (2,2)$; the residual $\vect{r}$ meets the line at a
right angle. The grey segments go to two other points of the line, which are further away.}
\label{fig:linalg:projection}
\end{figure}

\begin{worked}[label={we:linalg:proj2d}]{Projection in two dimensions}
Take $\vect{x} = (3, 1)$ and $\vect{u} = (1, 1)$, as in \cref{fig:linalg:projection}. Do this on
paper.
\begin{align*}
\vect{u}\T\vect{x} &= 1\times3 + 1\times1 = 3 + 1 = 4\\
\vect{u}\T\vect{u} &= 1\times1 + 1\times1 = 2\\
t^\star &= 4 / 2 = 2\\
\vect{p} &= 2\,(1, 1) = (2, 2)\\
\vect{r} &= (3, 1) - (2, 2) = (1, -1)\\
\vect{u}\T\vect{r} &= 1\times1 + 1\times(-1) = 0 .
\end{align*}
The residual is orthogonal to $\vect{u}$, as promised. Check that $t^\star = 2$ really is the
minimum by evaluating $f(t) = \norm{\vect{x}-t\vect{u}}^2$ at a few values of $t$:
\begin{center}
\begin{tabular}{cccc}
\toprule
$t$ & $\vect{x} - t\vect{u}$ & squares & $f(t)$\\
\midrule
1   & $(2, 0)$      & $4 + 0$         & 4\\
1.5 & $(1.5, -0.5)$ & $2.25 + 0.25$   & 2.5\\
2   & $(1, -1)$     & $1 + 1$         & 2\\
2.5 & $(0.5, -1.5)$ & $0.25 + 2.25$   & 2.5\\
3   & $(0, -2)$     & $0 + 4$         & 4\\
\bottomrule
\end{tabular}
\end{center}
The values fall to 2 at $t = 2$ and rise symmetrically on either side: a parabola with its
bottom at $t^\star$. Finally, Pythagoras: $\norm{\vect{x}}^2 = 9 + 1 = 10$, and
$\norm{\vect{p}}^2 + \norm{\vect{r}}^2 = (4 + 4) + (1 + 1) = 10$.
\end{worked}

\begin{worked}[label={we:linalg:proj3b}]{Projecting a pixel onto a reference spectrum}
A reference spectrum has the shape $\vect{u} = (0.1, 0.2, 0.3)$. A pixel reads
$\vect{x} = (0.12, 0.24, 0.30)$. How much of the pixel is ``the reference, scaled'', and what is
left over? Compute the two dot products:
\begin{align*}
\vect{u}\T\vect{x} &= 0.1\times0.12 + 0.2\times0.24 + 0.3\times0.30
  = 0.012 + 0.048 + 0.090 = 0.150\\
\vect{u}\T\vect{u} &= 0.01 + 0.04 + 0.09 = 0.14 .
\end{align*}
The best multiple is
\[
t^\star = \frac{0.150}{0.14} = 1.0714 ,
\]
so the projection is
\begin{align*}
\vect{p} = 1.0714 \times (0.1, 0.2, 0.3)
&= (0.10714,\ 0.21429,\ 0.32143)
\end{align*}
and the residual is
\begin{align*}
\vect{r} = \vect{x} - \vect{p}
&= (0.12 - 0.10714,\ 0.24 - 0.21429,\ 0.30 - 0.32143)\\
&= (0.01286,\ 0.02571,\ -0.02143).
\end{align*}
Check orthogonality:
$0.1\times0.01286 + 0.2\times0.02571 + 0.3\times(-0.02143) = 0.001286 + 0.005143 - 0.006429 = 0$.
The pixel is the reference at 107\% brightness, plus a small residual that is a little high in
bands 1 and 2 and a little low in band 3. Its squared length,
$0.01286^2 + 0.02571^2 + 0.02143^2 = 0.000165 + 0.000661 + 0.000459 = 0.001286$, is under 1\% of
$\norm{\vect{x}}^2 = 0.0144 + 0.0576 + 0.09 = 0.162$: the pixel is almost entirely the
reference shape.
\end{worked}

\begin{keypoint}
Projection onto a unit vector is one dot product, and what it leaves behind is at right angles
to the direction you projected on.
\end{keypoint}

\section{Matrices, the matrix-vector product and the transpose}\label{sec:linalg:matrices}

\begin{definition}[Matrix]\index{matrix}
A \emph{matrix} with $m$ rows and $n$ columns is a rectangular table of numbers; $A_{ij}$ is the
entry in row $i$, column $j$. A vector of length $n$ is a matrix with $n$ rows and one column.
\end{definition}

\subsection{Matrix times vector}

\begin{definition}[Matrix-vector product]\index{matrix-vector product}
For an $m \times n$ matrix $\mat{A}$ and a vector $\vect{x}$ of length $n$, the product
$\vect{y} = \mat{A}\vect{x}$ is the vector of length $m$ with
\begin{equation}\label{eq:linalg:matvec}
y_i = \sum_{j=1}^{n} A_{ij}\,x_j = (\text{row $i$ of } \mat{A})\T\vect{x}.
\end{equation}
\end{definition}

Read \cref{eq:linalg:matvec} in two ways. Row by row, each output entry is the dot product of
one row of $\mat{A}$ with $\vect{x}$, so each output is a weighted mixture of all the inputs with
its own weights. If the rows of $\mat{A}$ are unit vectors, the outputs are the coordinates of
$\vect{x}$ along those directions (\cref{sec:linalg:projection}); this is exactly how a
pixel is turned onto its principal directions later in the chapter.

\begin{worked}[label={we:linalg:matvec}]{A matrix times a vector}
Let
$\mat{A} = \begin{pmatrix} 1 & 2 \\ 3 & 4 \end{pmatrix}$ and
$\vect{v} = \begin{pmatrix} 1 \\ -1 \end{pmatrix}$.
\begin{align*}
y_1 &= (\text{row 1})\T\vect{v} = 1\times1 + 2\times(-1) = 1 - 2 = -1\\
y_2 &= (\text{row 2})\T\vect{v} = 3\times1 + 4\times(-1) = 3 - 4 = -1
\end{align*}
so $\mat{A}\vect{v} = (-1, -1)$.
\end{worked}

\subsection{Matrix times matrix}

\begin{definition}[Matrix product]\index{matrix product}
For an $m \times n$ matrix $\mat{A}$ and an $n \times p$ matrix $\mat{B}$, the product
$\mat{A}\mat{B}$ is the $m \times p$ matrix with
\[
(\mat{A}\mat{B})_{ik} = \sum_{j=1}^{n} A_{ij} B_{jk}
= (\text{row $i$ of } \mat{A})\T(\text{column $k$ of } \mat{B}).
\]
The inner sizes must agree ($n$ and $n$). Applying $\mat{A}\mat{B}$ to a vector is the same as
applying $\mat{B}$ first and then $\mat{A}$: $(\mat{A}\mat{B})\vect{x} = \mat{A}(\mat{B}\vect{x})$.
\end{definition}

\begin{definition}[Transpose]\index{transpose}
The \emph{transpose} $\mat{A}\T$ of an $m\times n$ matrix is the $n\times m$ matrix with the rows
and columns swapped: $(\mat{A}\T)_{ij} = A_{ji}$. A column vector transposed is a row vector,
which is why $\vect{a}\T\vect{b}$ (a row times a column) is the dot product. A matrix equal to its
own transpose, $\mat{S}\T = \mat{S}$, is \emph{symmetric}\index{symmetric matrix}.
\end{definition}

The rule you will use constantly is that transposing a product reverses the order:
\begin{equation}\label{eq:linalg:abt}
(\mat{A}\mat{B})\T = \mat{B}\T\mat{A}\T .
\end{equation}
The proof is one line of index juggling. Entry $(i,k)$ of the left side is entry $(k,i)$ of
$\mat{A}\mat{B}$, which is $\sum_j A_{kj}B_{ji}$. Entry $(i,k)$ of the right side is
$\sum_j (\mat{B}\T)_{ij}(\mat{A}\T)_{jk} = \sum_j B_{ji}A_{kj}$. The two sums contain the same
products, so they are equal. Order matters because matrix products do not commute in general:
$\mat{A}\mat{B} \neq \mat{B}\mat{A}$.

\begin{worked}[label={we:linalg:abt}]{The transpose of a product, on $2\times2$ matrices}
Let $\mat{A} = \begin{pmatrix} 1 & 2 \\ 3 & 4 \end{pmatrix}$ and
$\mat{B} = \begin{pmatrix} 0 & 1 \\ 5 & 2 \end{pmatrix}$. Do this on paper.
First $\mat{A}\mat{B}$, entry by entry (row of $\mat{A}$ with column of $\mat{B}$):
\begin{align*}
(1,1):&\quad 1\times0 + 2\times5 = 0 + 10 = 10\\
(1,2):&\quad 1\times1 + 2\times2 = 1 + 4 = 5\\
(2,1):&\quad 3\times0 + 4\times5 = 0 + 20 = 20\\
(2,2):&\quad 3\times1 + 4\times2 = 3 + 8 = 11
\end{align*}
so $\mat{A}\mat{B} = \begin{pmatrix} 10 & 5 \\ 20 & 11 \end{pmatrix}$ and
$(\mat{A}\mat{B})\T = \begin{pmatrix} 10 & 20 \\ 5 & 11 \end{pmatrix}$.
Now $\mat{B}\T = \begin{pmatrix} 0 & 5 \\ 1 & 2 \end{pmatrix}$ and
$\mat{A}\T = \begin{pmatrix} 1 & 3 \\ 2 & 4 \end{pmatrix}$, and
\begin{align*}
(1,1):&\quad 0\times1 + 5\times2 = 0 + 10 = 10\\
(1,2):&\quad 0\times3 + 5\times4 = 0 + 20 = 20\\
(2,1):&\quad 1\times1 + 2\times2 = 1 + 4 = 5\\
(2,2):&\quad 1\times3 + 2\times4 = 3 + 8 = 11
\end{align*}
so $\mat{B}\T\mat{A}\T = \begin{pmatrix} 10 & 20 \\ 5 & 11 \end{pmatrix}$, equal to
$(\mat{A}\mat{B})\T$. The wrong order gives something else:
$\mat{A}\T\mat{B}\T = \begin{pmatrix} 1\times0+3\times1 & 1\times5+3\times2 \\
2\times0+4\times1 & 2\times5+4\times2 \end{pmatrix} = \begin{pmatrix} 3 & 11 \\ 4 & 18
\end{pmatrix}$.
\end{worked}

\subsection{The identity and the inverse}

The \emph{identity matrix}\index{identity matrix} $\mat{I}$ has ones on the diagonal and zeros
elsewhere; $\mat{I}\vect{x} = \vect{x}$ for every $\vect{x}$. A square matrix $\mat{A}$ is
\emph{invertible} if there is a matrix $\mat{A}\inv$ with $\mat{A}\mat{A}\inv = \mat{A}\inv\mat{A}
= \mat{I}$. The inverse undoes the matrix: if $\vect{y} = \mat{A}\vect{x}$ then
$\vect{x} = \mat{A}\inv\vect{y}$.

\begin{definition}[Inverse of a $2\times2$ matrix]\index{inverse!of a 2x2 matrix}
For $\mat{A} = \begin{pmatrix} a & b \\ c & d \end{pmatrix}$ with \emph{determinant}
$\det\mat{A} = ad - bc \neq 0$,
\begin{equation}\label{eq:linalg:inv2}
\mat{A}\inv = \frac{1}{ad - bc}\begin{pmatrix} d & -b \\ -c & a \end{pmatrix}.
\end{equation}
\end{definition}

You do not need to memorise a derivation; you need to be able to check it. Multiply:
\[
\begin{pmatrix} a & b \\ c & d \end{pmatrix}\begin{pmatrix} d & -b \\ -c & a \end{pmatrix}
= \begin{pmatrix} ad - bc & -ab + ba \\ cd - dc & -cb + da \end{pmatrix}
= \begin{pmatrix} ad - bc & 0 \\ 0 & ad - bc \end{pmatrix}
= (ad - bc)\,\mat{I}.
\]
Dividing by $ad - bc$ gives $\mat{I}$. The recipe is: swap the diagonal, negate the
off-diagonal, divide by the determinant. When the determinant is zero there is no inverse: the
matrix squashes some direction flat and nothing can undo that. When the determinant is
\emph{nearly} zero, the inverse exists but contains huge numbers. \Cref{ch:rx} is in part a
chapter about that second case.\index{determinant}

\begin{worked}[label={we:linalg:inv}]{Inverting a $2\times2$ matrix by hand}
Take $\mat{S} = \begin{pmatrix} 2 & 3 \\ 3 & 10 \end{pmatrix}$, which is the covariance matrix
you will build in \cref{we:linalg:cov}.
\begin{align*}
\det\mat{S} &= 2\times10 - 3\times3 = 20 - 9 = 11\\
\mat{S}\inv &= \frac{1}{11}\begin{pmatrix} 10 & -3 \\ -3 & 2 \end{pmatrix}
= \begin{pmatrix} 0.9091 & -0.2727 \\ -0.2727 & 0.1818 \end{pmatrix}.
\end{align*}
Check the first column of $\mat{S}\mat{S}\inv$ (using the fractions, not the decimals):
$\tfrac{1}{11}(2\times10 + 3\times(-3)) = \tfrac{1}{11}(20 - 9) = 1$ and
$\tfrac{1}{11}(3\times10 + 10\times(-3)) = \tfrac{1}{11}(30 - 30) = 0$. Do the second column
yourself; you should get 0 and 1.
\end{worked}

\section{Mean, centring and the covariance matrix}\label{sec:linalg:covariance}

Now the statistics. Suppose a scene has $N$ valid pixels with spectra
$\vect{x}_1, \dots, \vect{x}_N$, each in $\R^C$.

\begin{definition}[Mean vector]\index{mean vector}
The \emph{mean} spectrum is the band-by-band average
\begin{equation}\label{eq:linalg:mean}
\muv = \frac{1}{N}\sum_{n=1}^{N}\vect{x}_n .
\end{equation}
\end{definition}

\begin{definition}[Centring]\index{centring}
\emph{Centring} subtracts the mean from every pixel: $\vect{d}_n = \vect{x}_n - \muv$. The
centred vectors (\emph{deviations}) always add to the zero vector, because
$\sum_n \vect{d}_n = \sum_n\vect{x}_n - N\muv = N\muv - N\muv = \vect{0}$.
\end{definition}

The \emph{variance}\index{variance} of one band is the average squared deviation in that band.
The \emph{covariance}\index{covariance} of two bands $i$ and $j$ is the average of the product of
their deviations, $d_{n,i}\,d_{n,j}$. When the two bands tend to be above their means together and
below together, most products are positive and the covariance is positive; when one tends to be
up while the other is down, it is negative; when they are unrelated, positive and negative
products cancel and it is near zero. The variance is the covariance of a band with itself.

All $C^2$ of these numbers fit into one matrix, and the neat way to build it is from
\emph{outer products}.

\begin{definition}[Outer product]\index{outer product}
For a column vector $\vect{d}$ of length $C$, the outer product $\vect{d}\vect{d}\T$ is the
$C \times C$ matrix whose entry $(i,j)$ is $d_i d_j$. (Compare the dot product $\vect{d}\T\vect{d}$,
a single number: the same symbols in the other order.)
\end{definition}

Entry $(i,j)$ of $\vect{d}_n\vect{d}_n\T$ is exactly the product $d_{n,i}d_{n,j}$ that the
covariance of bands $i$ and $j$ averages. So averaging the outer products averages every
pairwise product at once:

\begin{definition}[Covariance matrix]\index{covariance matrix}
The (sample) \emph{covariance matrix} of the pixels is
\begin{equation}\label{eq:linalg:cov}
\Cov = \frac{1}{N-1}\sum_{n=1}^{N} (\vect{x}_n - \muv)(\vect{x}_n - \muv)\T
     = \frac{1}{N-1}\sum_{n=1}^{N} \vect{d}_n\vect{d}_n\T .
\end{equation}
Its diagonal holds the band variances; entry $(i,j)$ off the diagonal holds the covariance of
bands $i$ and $j$. Because $d_i d_j = d_j d_i$, it is symmetric.
\end{definition}

Why $N-1$ and not $N$? The deviations are measured from the sample's own mean, which sits a
little closer to the sample than the true mean does, so the average squared deviation slightly
underestimates the true spread. Dividing by $N-1$ instead of $N$ corrects that on average. The
deeper reason, which \cref{ch:localrx} needs, is that the $N$ deviations add up to zero, so only
$N-1$ of them are free. With millions of pixels the difference between the two divisors is
negligible; with a handful it is not, and you must know which one the code uses.

\begin{sourcenote}[: which divisor the code uses]
\code{np.cov} divides by $N-1$, and every multi-band detector in this book computes its
covariance that way: the \code{spectral} library behind \code{spectral.rx} calls \code{np.cov}
on array input (\file{spectral/algorithms/algorithms.py}, function \code{mean_cov}, version 0.23.1),
\code{app/detectors/_local_background.py} line 68 divides by \code{counts_f - 1}, and
\code{app/utils/anomaly_detection/residual_mahalanobis.py} line 44 calls \code{np.cov}. The one
exception is the single-band thermal detector,
\code{app/detectors/thermal_grx_detector.py} line 100, which uses \code{vals.var()}, numpy's
default divisor $N$. On a scene of millions of pixels the two agree to many digits.
\end{sourcenote}

\begin{worked}[label={we:linalg:cov}]{A covariance matrix from four pixels and two bands}
Four pixels, two bands, reflectance in percent. Do every line on paper.
\begin{center}
\begin{tabular}{ccc}
\toprule
pixel & band 1 & band 2\\
\midrule
1 & 8  & 18\\
2 & 10 & 17\\
3 & 11 & 24\\
4 & 11 & 21\\
\bottomrule
\end{tabular}
\end{center}
\textbf{Step 1, the mean.}
\begin{align*}
\mu_1 &= (8 + 10 + 11 + 11)/4 = 40/4 = 10\\
\mu_2 &= (18 + 17 + 24 + 21)/4 = 80/4 = 20 .
\end{align*}
\textbf{Step 2, centre.} Subtract $(10, 20)$ from each pixel:
\begin{align*}
\vect{d}_1 &= (8-10,\ 18-20) = (-2, -2)\\
\vect{d}_2 &= (10-10,\ 17-20) = (0, -3)\\
\vect{d}_3 &= (11-10,\ 24-20) = (1, 4)\\
\vect{d}_4 &= (11-10,\ 21-20) = (1, 1).
\end{align*}
Check that they add to zero: band 1, $-2 + 0 + 1 + 1 = 0$; band 2, $-2 - 3 + 4 + 1 = 0$.

\textbf{Step 3, the outer products.} Entry $(i,j)$ is $d_i d_j$:
\begin{align*}
\vect{d}_1\vect{d}_1\T &= \begin{pmatrix} (-2)(-2) & (-2)(-2) \\ (-2)(-2) & (-2)(-2)\end{pmatrix}
  = \begin{pmatrix} 4 & 4 \\ 4 & 4 \end{pmatrix}\\
\vect{d}_2\vect{d}_2\T &= \begin{pmatrix} 0\times0 & 0\times(-3) \\ (-3)\times0 & (-3)(-3)\end{pmatrix}
  = \begin{pmatrix} 0 & 0 \\ 0 & 9 \end{pmatrix}\\
\vect{d}_3\vect{d}_3\T &= \begin{pmatrix} 1\times1 & 1\times4 \\ 4\times1 & 4\times4\end{pmatrix}
  = \begin{pmatrix} 1 & 4 \\ 4 & 16 \end{pmatrix}\\
\vect{d}_4\vect{d}_4\T &= \begin{pmatrix} 1\times1 & 1\times1 \\ 1\times1 & 1\times1\end{pmatrix}
  = \begin{pmatrix} 1 & 1 \\ 1 & 1 \end{pmatrix}.
\end{align*}
\textbf{Step 4, add them} entry by entry:
\begin{align*}
(1,1):&\quad 4 + 0 + 1 + 1 = 6\\
(1,2) = (2,1):&\quad 4 + 0 + 4 + 1 = 9\\
(2,2):&\quad 4 + 9 + 16 + 1 = 30 .
\end{align*}
\textbf{Step 5, divide by $N - 1 = 3$:}
\[
\Cov = \frac{1}{3}\begin{pmatrix} 6 & 9 \\ 9 & 30 \end{pmatrix}
     = \begin{pmatrix} 2 & 3 \\ 3 & 10 \end{pmatrix}.
\]
Read it: band 1 has variance 2 (spread $\sqrt2 = 1.41$), band 2 has variance 10 (spread
$\sqrt{10} = 3.16$), and their covariance is $+3$: they tend to rise and fall together. The
\emph{correlation}\index{correlation}, the covariance divided by both spreads, is
$3/(\sqrt2\times\sqrt{10}) = 3/4.472 = 0.67$. Had you divided by $N = 4$ you would have got
$\begin{pmatrix} 1.5 & 2.25 \\ 2.25 & 7.5\end{pmatrix}$, smaller by the factor $3/4$.
\end{worked}

In matrix form the whole calculation is one line. Stack the centred pixels as the rows of an
$N \times C$ matrix $\mat{D}$. Then $\mat{D}\T\mat{D}$ is $C\times C$, and its entry $(i,j)$ is
$\sum_n d_{n,i}d_{n,j}$: exactly the sum of outer products. So
\begin{equation}\label{eq:linalg:covmat}
\Cov = \frac{1}{N-1}\,\mat{D}\T\mat{D}.
\end{equation}
This is the line \code{cov = dX.transpose(-1, -2) @ dX} in
\code{app/detectors/_local_background.py}, done for a whole batch of pixels at once.

\subsection{How a linear map changes the covariance}

Almost everything in the rest of Part II transforms the pixels by a matrix and asks what
happens to their mean and covariance. The answer is short and worth deriving once. Let
$\vect{y}_n = \mat{A}\vect{x}_n$ for a fixed matrix $\mat{A}$.

The mean of the $\vect{y}$'s is $\mat{A}\muv$, because a matrix product distributes over the sum:
$\frac1N\sum_n\mat{A}\vect{x}_n = \mat{A}\,\frac1N\sum_n\vect{x}_n$. So the deviations are
$\vect{y}_n - \mat{A}\muv = \mat{A}\vect{d}_n$, and
\begin{align}
\Cov_y &= \frac{1}{N-1}\sum_n (\mat{A}\vect{d}_n)(\mat{A}\vect{d}_n)\T
  && \text{definition, \cref{eq:linalg:cov}} \notag\\
&= \frac{1}{N-1}\sum_n \mat{A}\vect{d}_n\vect{d}_n\T\mat{A}\T
  && \text{transpose of a product, \cref{eq:linalg:abt}} \notag\\
&= \mat{A}\left(\frac{1}{N-1}\sum_n \vect{d}_n\vect{d}_n\T\right)\mat{A}\T
  && \text{$\mat{A}$ is the same in every term; factor it out} \notag\\
&= \mat{A}\,\Cov\,\mat{A}\T . \label{eq:linalg:covtransform}
\end{align}

\begin{keypoint}
If $\vect{y} = \mat{A}\vect{x}$, then $\Cov_y = \mat{A}\Cov\mat{A}\T$.
\end{keypoint}

A special case is a single direction. Take $\mat{A}$ to be one row, $\vect{u}\T$ with $\vect{u}$ a
unit vector. Then $y_n = \vect{u}\T\vect{x}_n$ is the coordinate of each pixel along $\vect{u}$,
and its variance is the $1\times1$ matrix
\begin{equation}\label{eq:linalg:dirvar}
\operatorname{Var}(\vect{u}\T\vect{x}) = \vect{u}\T\Cov\,\vect{u}.
\end{equation}
The covariance matrix therefore knows the spread of the cloud in \emph{every} direction, not only
along the band axes. For the matrix of \cref{we:linalg:cov} and $\vect{u} = (1, 0)$ you get
$\vect{u}\T\Cov\vect{u} = \Sigma_{11} = 2$, the variance of band 1. For the diagonal
$\vect{u} = (1, 1)/\sqrt2$:
\[
\vect{u}\T\Cov\vect{u} = \tfrac12\,(1, 1)\begin{pmatrix} 2 & 3\\ 3 & 10\end{pmatrix}
\begin{pmatrix}1\\1\end{pmatrix} = \tfrac12\,(1, 1)\begin{pmatrix} 5\\ 13\end{pmatrix}
= \tfrac12\,(5 + 13) = 9 .
\]
Along the diagonal the cloud spreads more than along either axis. Which direction spreads the
most? That is the question eigenvectors answer.

\section{Eigenvalues and eigenvectors}\label{sec:linalg:eigen}

\begin{definition}[Eigenvector and eigenvalue]\index{eigenvector}\index{eigenvalue}
A non-zero vector $\vect{u}$ is an \emph{eigenvector} of a square matrix $\mat{S}$, with
\emph{eigenvalue} $\lambda$, if
\begin{equation}\label{eq:linalg:eig}
\mat{S}\vect{u} = \lambda\vect{u}.
\end{equation}
The matrix does not turn an eigenvector; it only stretches it by $\lambda$.
\end{definition}

Most vectors are turned by a matrix. For $\mat{S} = \begin{pmatrix} 2 & 3\\ 3 & 10\end{pmatrix}$,
the vector $(1, 0)$ goes to $(2, 3)$, a different direction. We want the special directions that
are only stretched.

\subsection{Finding them by hand: the characteristic polynomial}

Rewrite \cref{eq:linalg:eig} as $\mat{S}\vect{u} - \lambda\vect{u} = \vect{0}$, that is
$(\mat{S} - \lambda\mat{I})\vect{u} = \vect{0}$. A non-zero $\vect{u}$ is sent to zero, so the
matrix $\mat{S} - \lambda\mat{I}$ squashes a direction flat, and (\cref{eq:linalg:inv2}) its
determinant must be zero. For a $2\times2$ matrix
$\mat{S} = \begin{pmatrix} a & b\\ b & c\end{pmatrix}$ (symmetric: both off-diagonal entries are
$b$):
\begin{align}
\det(\mat{S} - \lambda\mat{I}) &= (a - \lambda)(c - \lambda) - b^2 \notag\\
&= \lambda^2 - (a + c)\lambda + (ac - b^2) = 0 .
  \label{eq:linalg:charpoly}
\end{align}
This quadratic is the \emph{characteristic polynomial}\index{characteristic polynomial}. Its
coefficients are the \emph{trace} $a + c$ (the sum of the diagonal) and the determinant
$ac - b^2$. Its two roots are the eigenvalues:
\begin{equation}\label{eq:linalg:eig2}
\lambda_{1,2} = \frac{(a + c) \pm \sqrt{(a - c)^2 + 4b^2}}{2}.
\end{equation}
(Expand $(a+c)^2 - 4(ac - b^2) = a^2 - 2ac + c^2 + 4b^2 = (a-c)^2 + 4b^2$ to get the square root.)
The quantity under the root is a sum of squares, so it is never negative: a symmetric
$2\times2$ matrix always has two real eigenvalues. For each eigenvalue, the first row of
$(\mat{S} - \lambda\mat{I})\vect{u} = \vect{0}$, namely $(a - \lambda)u_1 + b u_2 = 0$, gives the
direction.

\begin{worked}[label={we:linalg:eig}]{Eigenvalues and eigenvectors of the covariance matrix}
For $\Cov = \begin{pmatrix} 2 & 3\\ 3 & 10\end{pmatrix}$: trace $2 + 10 = 12$, determinant
$2\times10 - 3\times3 = 11$. The characteristic polynomial is
\[
\lambda^2 - 12\lambda + 11 = 0 .
\]
It factorises: $(\lambda - 11)(\lambda - 1) = \lambda^2 - 11\lambda - \lambda + 11 = \lambda^2 -
12\lambda + 11$. So $\lambda_1 = 11$ and $\lambda_2 = 1$. (The formula agrees:
$\sqrt{(2-10)^2 + 4\times9} = \sqrt{64 + 36} = 10$, and $(12 \pm 10)/2 = 11$ or $1$.)

\textbf{Eigenvector for $\lambda_1 = 11$.} The first row gives
$(2 - 11)u_1 + 3u_2 = 0$, so $-9u_1 + 3u_2 = 0$, so $u_2 = 3u_1$. Take $(1, 3)$, length
$\sqrt{1 + 9} = \sqrt{10}$, and divide: $\vect{u}_1 = (1, 3)/\sqrt{10} = (0.3162, 0.9487)$.
Check without the scaling: $\Cov\,(1, 3) = (2\times1 + 3\times3,\ 3\times1 + 10\times3) = (11, 33)
= 11\times(1, 3)$.

\textbf{Eigenvector for $\lambda_2 = 1$.} $(2 - 1)u_1 + 3u_2 = 0$, so $u_1 = -3u_2$. Take
$(3, -1)$ and divide by $\sqrt{10}$: $\vect{u}_2 = (3, -1)/\sqrt{10} = (0.9487, -0.3162)$.
Check: $\Cov\,(3, -1) = (6 - 3,\ 9 - 10) = (3, -1) = 1\times(3, -1)$.

The two eigenvectors are orthogonal: $(1, 3)\T(3, -1) = 3 - 3 = 0$.
\end{worked}

The orthogonality is not luck. For any symmetric $\mat{S}$, take eigenvectors $\vect{u}$ and
$\vect{v}$ with different eigenvalues $\lambda \ne \nu$. Compute $\vect{v}\T\mat{S}\vect{u}$ two
ways: as $\vect{v}\T(\lambda\vect{u}) = \lambda\,\vect{v}\T\vect{u}$, and, using
$\mat{S}\T = \mat{S}$ and \cref{eq:linalg:abt}, as $(\mat{S}\vect{v})\T\vect{u} =
\nu\,\vect{v}\T\vect{u}$. So $(\lambda - \nu)\,\vect{v}\T\vect{u} = 0$, and since
$\lambda \ne \nu$ the dot product must be zero.

\subsection{The spectral theorem}

\begin{definition}[Spectral theorem for symmetric matrices]\index{spectral theorem}
Every real symmetric $C\times C$ matrix $\mat{S}$ has $C$ real eigenvalues
$\lambda_1 \ge \lambda_2 \ge \cdots \ge \lambda_C$ and a matching set of $C$ orthogonal unit
eigenvectors $\vect{u}_1, \dots, \vect{u}_C$. Put the eigenvectors as the columns of
$\mat{U}$ and the eigenvalues on the diagonal of $\mat{\Lambda} = \diag(\lambda_1, \dots,
\lambda_C)$. Then
\begin{equation}\label{eq:linalg:spectral}
\mat{U}\T\mat{U} = \mat{I}
\qquad\text{and}\qquad
\mat{S} = \mat{U}\mat{\Lambda}\mat{U}\T = \sum_{k=1}^{C}\lambda_k\,\vect{u}_k\vect{u}_k\T .
\end{equation}
A square matrix with $\mat{U}\T\mat{U} = \mat{I}$ is \emph{orthogonal}\index{orthogonal matrix}:
it only rotates (or reflects), never stretches, and its inverse is its transpose.
\end{definition}

We prove the two-by-two case by multiplication in the worked example below and take the general
case on trust; it is proved in every linear algebra text. The decomposition is what
\code{np.linalg.eigh} computes; the \code{h} stands for Hermitian, the complex-number word for
symmetric. It returns the eigenvalues in \emph{ascending} order, which is why every piece of code
in this repository that uses it reverses the order with \code{np.argsort(evals)[::-1]}.

A covariance matrix has one more property: its eigenvalues are never negative. By
\cref{eq:linalg:dirvar}, $\lambda_k = \vect{u}_k\T\Cov\vect{u}_k$ is the variance of the pixels
along $\vect{u}_k$, and a variance is an average of squares. A matrix with all eigenvalues
$\ge 0$ is \emph{positive semi-definite}\index{positive semi-definite}; if all are $>0$ it is
\emph{positive definite} and invertible.

\begin{worked}[label={we:linalg:spectral}]{Rebuilding the covariance from its eigenvectors}
With $\vect{u}_1 = (1, 3)/\sqrt{10}$, $\vect{u}_2 = (3, -1)/\sqrt{10}$, $\lambda_1 = 11$,
$\lambda_2 = 1$, use the sum form of \cref{eq:linalg:spectral}. The outer products are
\[
\vect{u}_1\vect{u}_1\T = \frac{1}{10}\begin{pmatrix} 1 & 3\\ 3 & 9\end{pmatrix},
\qquad
\vect{u}_2\vect{u}_2\T = \frac{1}{10}\begin{pmatrix} 9 & -3\\ -3 & 1\end{pmatrix}.
\]
Weight them by the eigenvalues and add:
\[
11\times\frac{1}{10}\begin{pmatrix} 1 & 3\\ 3 & 9\end{pmatrix}
+ 1\times\frac{1}{10}\begin{pmatrix} 9 & -3\\ -3 & 1\end{pmatrix}
= \frac{1}{10}\begin{pmatrix} 11 + 9 & 33 - 3\\ 33 - 3 & 99 + 1\end{pmatrix}
= \frac{1}{10}\begin{pmatrix} 20 & 30\\ 30 & 100\end{pmatrix}
= \begin{pmatrix} 2 & 3\\ 3 & 10\end{pmatrix}.
\]
The covariance is back. Read the sum form as a statement about the cloud: it is the sum of a
spread of 11 along $\vect{u}_1$ and a spread of 1 along $\vect{u}_2$, and nothing else.
\end{worked}

\section{Principal component analysis}\label{sec:linalg:pca}

\subsection{PCA as variance maximisation}

\emph{Principal component analysis}\index{principal component analysis} (PCA) asks: along which
unit direction $\vect{u}$ do the pixels spread the most? By \cref{eq:linalg:dirvar} the spread
along $\vect{u}$ is $\vect{u}\T\Cov\vect{u}$, so the question is the constrained maximisation
\begin{equation}\label{eq:linalg:pcaproblem}
\max_{\vect{u}}\ \vect{u}\T\Cov\vect{u}
\qquad\text{subject to}\qquad \vect{u}\T\vect{u} = 1 .
\end{equation}
The constraint matters: without it you could make $\vect{u}\T\Cov\vect{u}$ as large as you like by
making $\vect{u}$ long.

The standard tool for maximising under a constraint is a \emph{Lagrange
multiplier}\index{Lagrange multiplier}. The idea in one sentence: at the best point you cannot
increase the objective by moving along the constraint, so the objective's gradient must point
straight out of the constraint surface, which means it is a multiple $\lambda$ of the constraint's
gradient. Formally, build
\[
L(\vect{u}, \lambda) = \vect{u}\T\Cov\vect{u} - \lambda\,(\vect{u}\T\vect{u} - 1)
\]
and set its derivatives with respect to every entry of $\vect{u}$ to zero.

We need the derivative of a quadratic form. Work it out on two bands first. With
$\Cov = \begin{pmatrix} a & b\\ b & c\end{pmatrix}$ and $\vect{u} = (u_1, u_2)$, multiplying out
gives
\[
\vect{u}\T\Cov\vect{u} = a u_1^2 + 2b\,u_1u_2 + c u_2^2 .
\]
Differentiate with respect to $u_1$ (treat $u_2$ as a constant) and with respect to $u_2$:
\begin{align*}
\frac{\partial}{\partial u_1}\left(a u_1^2 + 2bu_1u_2 + cu_2^2\right) &= 2au_1 + 2bu_2,\\
\frac{\partial}{\partial u_2}\left(a u_1^2 + 2bu_1u_2 + cu_2^2\right) &= 2bu_1 + 2cu_2 .
\end{align*}
Stacked, these are $2\begin{pmatrix} a & b\\ b & c\end{pmatrix}\begin{pmatrix}u_1\\u_2\end{pmatrix}
= 2\Cov\vect{u}$. The same holds in any number of dimensions for a symmetric matrix: the gradient
of $\vect{u}\T\Cov\vect{u}$ is $2\Cov\vect{u}$. With $\Cov = \mat{I}$ it says the gradient of
$\vect{u}\T\vect{u}$ is $2\vect{u}$.

Now set the gradient of $L$ to zero:
\begin{align}
2\Cov\vect{u} - 2\lambda\vect{u} &= \vect{0}
  && \text{gradient of each term} \notag\\
\Cov\vect{u} &= \lambda\vect{u}
  && \text{add $2\lambda\vect{u}$, divide by 2.} \label{eq:linalg:pcaeig}
\end{align}
The best direction is an eigenvector of $\Cov$. Which one? Multiply \cref{eq:linalg:pcaeig} on the
left by $\vect{u}\T$ and use $\vect{u}\T\vect{u} = 1$:
\[
\vect{u}\T\Cov\vect{u} = \lambda\,\vect{u}\T\vect{u} = \lambda .
\]
The variance along an eigenvector is its eigenvalue, so the maximum is reached at the eigenvector
with the \emph{largest} eigenvalue. That eigenvector is the \emph{first principal
direction}\index{principal direction}, and the coordinate of a pixel along it,
$\vect{u}_1\T(\vect{x} - \muv)$, is its \emph{first principal component}.

The second principal direction maximises the variance among unit directions orthogonal to the
first. Repeating the argument with a second multiplier for the constraint
$\vect{u}\T\vect{u}_1 = 0$ gives the eigenvector with the second-largest eigenvalue, and so on: the
principal directions are the eigenvectors of $\Cov$, in order of decreasing eigenvalue.

\begin{worked}[label={we:linalg:pcacheck}]{Checking that the first principal direction wins}
For the covariance of \cref{we:linalg:cov}, compute the spread
$\vect{u}\T\Cov\vect{u} = 2u_1^2 + 6u_1u_2 + 10u_2^2$ along unit directions at a few angles
$\phi$ from the band-1 axis, $\vect{u} = (\cos\phi, \sin\phi)$:
\begin{center}
\begin{tabular}{cccc}
\toprule
$\phi$ & $\vect{u}$ & $2u_1^2 + 6u_1u_2 + 10u_2^2$ & spread\\
\midrule
$0^\circ$  & $(1, 0)$          & $2 + 0 + 0$                 & 2.0\\
$45^\circ$ & $(0.7071, 0.7071)$ & $1 + 3 + 5$                & 9.0\\
$71.57^\circ$ & $(0.3162, 0.9487)$ & $0.2 + 1.8 + 9$       & 11.0\\
$90^\circ$ & $(0, 1)$          & $0 + 0 + 10$                & 10.0\\
$135^\circ$ & $(-0.7071, 0.7071)$ & $1 - 3 + 5$              & 3.0\\
\bottomrule
\end{tabular}
\end{center}
The line for $71.57^\circ$ is $\vect{u}_1 = (1,3)/\sqrt{10}$ (its angle is $\arctan 3$): there
$2\times0.1 = 0.2$, $6\times0.3 = 1.8$ and $10\times0.9 = 9$, using $u_1^2 = 0.1$,
$u_1u_2 = 0.3$, $u_2^2 = 0.9$. No direction beats 11, the largest eigenvalue, and the smallest
possible spread is the other eigenvalue, 1, at $\vect{u}_2$ (angle $-18.43^\circ$).
\end{worked}

\subsection{Explained variance}

Because $\mat{U}$ only rotates, the total variance is the same in the rotated coordinates as in
the bands: it is the trace, $\tr\Cov = \sum_i \Sigma_{ii} = \sum_k\lambda_k$. (For the example,
$2 + 10 = 12 = 11 + 1$.) The share carried by component $k$ is its \emph{explained variance
ratio}\index{explained variance ratio}
\begin{equation}\label{eq:linalg:evr}
\rho_k = \frac{\lambda_k}{\lambda_1 + \lambda_2 + \cdots + \lambda_C}.
\end{equation}
In the example $\rho_1 = 11/12 = 0.917$ and $\rho_2 = 1/12 = 0.083$: one direction carries 92\%
of the spread. Real hyperspectral scenes are far more extreme, because neighbouring bands of a
165-band spectrum are nearly copies of each other; a few components carry almost all the
variance and the remaining hundred-odd carry almost none. Keep that sentence in mind: it is the
whole of \cref{ch:rx}'s numerical problem.

\begin{worked}[label={we:linalg:pcacoords}]{Turning the four pixels onto their principal directions}
The principal components of each centred pixel are its dot products with $\vect{u}_1$ and
$\vect{u}_2$. Working with $\sqrt{10}\,\vect{u}_1 = (1, 3)$ and $\sqrt{10}\,\vect{u}_2 = (3, -1)$
keeps integers; divide by $\sqrt{10}$ at the end.
\begin{center}
\begin{tabular}{cccc}
\toprule
pixel & $\vect{d}$ & $(1,3)\T\vect{d}$ & $(3,-1)\T\vect{d}$\\
\midrule
1 & $(-2, -2)$ & $-2 - 6 = -8$  & $-6 + 2 = -4$\\
2 & $(0, -3)$  & $0 - 9 = -9$   & $0 + 3 = 3$\\
3 & $(1, 4)$   & $1 + 12 = 13$  & $3 - 4 = -1$\\
4 & $(1, 1)$   & $1 + 3 = 4$    & $3 - 1 = 2$\\
\bottomrule
\end{tabular}
\end{center}
Variance of the first component (the mean is $(-8-9+13+4)/\sqrt{10} = 0$):
\[
\frac{1}{3}\cdot\frac{(-8)^2 + (-9)^2 + 13^2 + 4^2}{10}
= \frac{1}{3}\cdot\frac{64 + 81 + 169 + 16}{10} = \frac{1}{3}\cdot\frac{330}{10} = 11 .
\]
Variance of the second:
$\frac{1}{3}\cdot\frac{16 + 9 + 1 + 4}{10} = \frac13\cdot\frac{30}{10} = 1$.
Covariance between them:
$\frac{1}{3}\cdot\frac{(-8)(-4) + (-9)(3) + (13)(-1) + (4)(2)}{10}
= \frac13\cdot\frac{32 - 27 - 13 + 8}{10} = 0$.
In the rotated coordinates the covariance is $\diag(11, 1) = \mat{\Lambda}$: the components are
uncorrelated, with variances equal to the eigenvalues. That is \cref{eq:linalg:covtransform}
with $\mat{A} = \mat{U}\T$: $\mat{U}\T\Cov\mat{U} = \mat{U}\T\mat{U}\mat{\Lambda}\mat{U}\T\mat{U}
= \mat{\Lambda}$.
\end{worked}

\section{Whitening}\label{sec:linalg:whitening}

After PCA the cloud's axes line up with the coordinate axes, but it is still stretched: spread
11 one way, 1 the other. \emph{Whitening}\index{whitening} also divides each rotated coordinate by
its own standard deviation, so the cloud becomes round with spread 1 in every direction. The
name comes from signal processing: ``white'' noise has equal power at every frequency, and a
whitened cloud has equal variance in every direction.

\begin{definition}[Whitening matrix]
With $\Cov = \mat{U}\mat{\Lambda}\mat{U}\T$ and all $\lambda_k > 0$, let
\[
\mat{\Lambda}^{-1/2} = \diag(1/\sqrt{\lambda_1}, \dots, 1/\sqrt{\lambda_C}).
\]
The \emph{whitening matrix} is
\begin{equation}\label{eq:linalg:whiten}
\mat{W} = \mat{\Lambda}^{-1/2}\mat{U}\T ,
\qquad
\vect{z} = \mat{W}(\vect{x} - \muv).
\end{equation}
Row $k$ of $\mat{W}$ is $\vect{u}_k\T/\sqrt{\lambda_k}$: ``take the coordinate along
$\vect{u}_k$, then divide by that direction's standard deviation''.
\end{definition}

Why the square root? The eigenvalue $\lambda_k$ is a variance, a squared spread. The coordinate
$\vect{u}_k\T\vect{d}$ is in the data's own units, so it must be divided by a spread in the same
units: $\sqrt{\lambda_k}$.

\paragraph{Claim.} If $\vect{z} = \mat{W}(\vect{x} - \muv)$ with $\mat{W}$ as in
\cref{eq:linalg:whiten}, then $\Cov_z = \mat{I}$: the whitened pixels have identity covariance.

\noindent\textit{Proof.} By \cref{eq:linalg:covtransform} with $\mat{A} = \mat{W}$ (subtracting
$\muv$ does not change a covariance):
\begin{align*}
\Cov_z &= \mat{W}\Cov\mat{W}\T\\
&= \mat{\Lambda}^{-1/2}\mat{U}\T\;\mat{U}\mat{\Lambda}\mat{U}\T\;
   (\mat{\Lambda}^{-1/2}\mat{U}\T)\T
  && \text{substitute $\mat{W}$ and the spectral theorem}\\
&= \mat{\Lambda}^{-1/2}\mat{U}\T\mat{U}\mat{\Lambda}\mat{U}\T\mat{U}\mat{\Lambda}^{-1/2}
  && \text{\cref{eq:linalg:abt}; a diagonal matrix is its own transpose}\\
&= \mat{\Lambda}^{-1/2}\,\mat{I}\,\mat{\Lambda}\,\mat{I}\,\mat{\Lambda}^{-1/2}
  && \mat{U}\T\mat{U} = \mat{I}\\
&= \mat{\Lambda}^{-1/2}\mat{\Lambda}\mat{\Lambda}^{-1/2}
  && \text{the identity changes nothing}\\
&= \mat{I}
  && \text{diagonal entries: } \tfrac{1}{\sqrt{\lambda_k}}\,\lambda_k\,\tfrac{1}{\sqrt{\lambda_k}} = 1 .
\end{align*}
\hfill$\square$

\begin{worked}[label={we:linalg:whiten}]{Whitening the four pixels}
For the running example,
\[
\mat{W} = \begin{pmatrix} 1/\sqrt{11} & 0\\ 0 & 1\end{pmatrix}
\frac{1}{\sqrt{10}}\begin{pmatrix} 1 & 3\\ 3 & -1\end{pmatrix}
= \begin{pmatrix} 1/\sqrt{110} & 3/\sqrt{110}\\ 3/\sqrt{10} & -1/\sqrt{10}\end{pmatrix}
= \begin{pmatrix} 0.0953 & 0.2860\\ 0.9487 & -0.3162\end{pmatrix}.
\]
The whitened pixels are the principal components of \cref{we:linalg:pcacoords} divided by
$\sqrt{\lambda_k}$: the first by $\sqrt{11} = 3.317$, the second by $\sqrt1 = 1$. Pixel 1, for
example, had components $-8/\sqrt{10} = -2.530$ and $-4/\sqrt{10} = -1.265$, so
\[
\vect{z}_1 = (-2.530/3.317,\ -1.265/1) = (-0.763,\ -1.265).
\]
All four:
\begin{center}
\begin{tabular}{ccc}
\toprule
pixel & $z_1 = (1,3)\T\vect{d}/\sqrt{110}$ & $z_2 = (3,-1)\T\vect{d}/\sqrt{10}$\\
\midrule
1 & $-8/10.488 = -0.763$ & $-4/3.162 = -1.265$\\
2 & $-9/10.488 = -0.858$ & $3/3.162 = 0.949$\\
3 & $13/10.488 = 1.240$  & $-1/3.162 = -0.316$\\
4 & $4/10.488 = 0.381$   & $2/3.162 = 0.632$\\
\bottomrule
\end{tabular}
\end{center}
Variance of $z_1$: $\frac13(64 + 81 + 169 + 16)/110 = \frac13\cdot 330/110 = \frac13\cdot3 = 1$.
Variance of $z_2$: $\frac13(16 + 9 + 1 + 4)/10 = \frac13\cdot3 = 1$. Their covariance is
$\frac{1}{3}\cdot\frac{0}{\sqrt{1100}} = 0$ from the same products as before. The whitened
covariance is $\mat{I}$, as the proof said it must be.
\end{worked}

\begin{figure}[htbp]
\centering
\begin{tikzpicture}[>=Stealth]
  % left panel: raw cloud, scale 0.4 cm per unit
  \begin{scope}[scale=0.4]
    \draw[->,gray] (-5,0) -- (5,0) node[right,font=\scriptsize] {$d_1$};
    \draw[->,gray] (0,-7.5) -- (0,8) node[above,font=\scriptsize] {$d_2$};
    \draw[spat,thick,rotate=71.565] (0,0) ellipse ({sqrt(11)} and 1);
    \draw[spat,dashed,rotate=71.565] (0,0) ellipse ({2*sqrt(11)} and 2);
    \foreach \p in {(-0.01,-0.44),(-0.77,-1.98),(0.22,1.80),(-2.27,-3.80),(1.94,1.81),(0.61,1.97),(2.65,7.25),(2.00,3.69),(0.87,1.61),(-1.42,-0.97),(-1.16,0.84),(0.36,2.20),(-1.04,-2.81),(1.65,2.05),(0.27,1.51),(1.26,3.11),(-0.35,0.67),(-0.56,-0.56),(0.81,2.02),(-2.24,-1.99),(-0.12,4.44),(-2.15,-3.99),(0.02,-0.93),(-3.39,-3.97),(-3.14,-5.25),(0.77,-1.51),(-0.62,0.22),(0.98,4.18),(-1.11,-1.13),(1.87,3.52),(1.57,1.00),(-1.35,-5.63),(-0.10,-1.12),(-2.91,-4.56),(1.43,-2.69),(-0.70,0.05),(-0.47,-3.35),(-0.66,-6.15),(-0.19,2.38),(-1.95,0.48),(0.40,-0.79),(1.43,0.84),(0.74,4.22),(-1.84,-1.42),(-0.81,2.91),(-1.19,1.40),(1.15,1.13),(-1.24,-0.46),(0.39,-0.69),(2.74,3.78),(-1.47,-1.27),(1.19,0.18),(0.30,-3.15),(0.16,2.44),(1.99,7.20),(0.51,1.86),(-0.11,0.30),(-0.01,-2.42),(-2.18,-4.24),(-1.99,-2.29)}
      \fill[maskgrey] \p circle (0.12);
    \foreach \p in {(-2,-2),(0,-3),(1,4),(1,1)}
      \fill[spec] \p circle (0.22);
    \draw[->,very thick,loss] (0,0) -- ({sqrt(11)*0.3162},{sqrt(11)*0.9487})
      node[above right,font=\scriptsize] {$\sqrt{\lambda_1}\,\vect{u}_1$};
    \draw[->,very thick,attn] (0,0) -- (0.9487,-0.3162)
      node[below right,font=\scriptsize] {$\sqrt{\lambda_2}\,\vect{u}_2$};
    \node[font=\small] at (0,-9.2) {(a) centred pixels};
  \end{scope}
  % right panel: whitened cloud, scale 1.25 cm per unit
  \begin{scope}[xshift=7.6cm,scale=1.25]
    \draw[->,gray] (-2.6,0) -- (2.7,0) node[right,font=\scriptsize] {$z_1$};
    \draw[->,gray] (0,-2.3) -- (0,2.5) node[above,font=\scriptsize] {$z_2$};
    \draw[spat,thick] (0,0) circle (1);
    \draw[spat,dashed] (0,0) circle (2);
    \foreach \p in {(-0.13,0.13),(-0.64,-0.10),(0.54,-0.36),(-1.30,-0.95),(0.70,1.27),(0.62,-0.04),(2.33,0.22),(1.25,0.73),(0.54,0.32),(-0.41,-1.04),(0.13,-1.37),(0.67,-0.35),(-0.90,-0.09),(0.74,0.92),(0.46,-0.22),(1.01,0.21),(0.16,-0.54),(-0.21,-0.36),(0.65,0.13),(-0.78,-1.49),(1.26,-1.51),(-1.35,-0.78),(-0.26,0.31),(-1.46,-1.96),(-1.80,-1.32),(-0.36,1.21),(0.00,-0.66),(1.29,-0.40),(-0.43,-0.70),(1.18,0.66),(0.44,1.17),(-1.74,0.50),(-0.33,0.26),(-1.58,-1.32),(-0.63,2.20),(-0.05,-0.68),(-1.00,0.62),(-1.82,1.32),(0.66,-0.94),(-0.05,-2.00),(-0.19,0.63),(0.38,1.09),(1.28,-0.63),(-0.58,-1.29),(0.75,-1.69),(0.29,-1.57),(0.43,0.74),(-0.25,-1.03),(-0.16,0.59),(1.34,1.40),(-0.50,-0.99),(0.16,1.07),(-0.87,1.28),(0.71,-0.62),(2.25,-0.39),(0.58,-0.11),(0.08,-0.20),(-0.69,0.76),(-1.42,-0.73),(-0.84,-1.16)}
      \fill[maskgrey] \p circle (0.04);
    \foreach \p in {(-0.763,-1.265),(-0.858,0.949),(1.240,-0.316),(0.381,0.632)}
      \fill[spec] \p circle (0.07);
    \draw[->,very thick,loss] (0,0) -- (1,0);
    \draw[->,very thick,attn] (0,0) -- (0,1);
    \node[font=\small] at (0,-2.85) {(b) whitened, $\vect{z} = \mat{W}\vect{d}$};
  \end{scope}
\end{tikzpicture}
\caption{Before and after whitening. (a) Sixty points drawn from a Gaussian with the covariance
$\begin{psmallmatrix} 2 & 3\\ 3 & 10\end{psmallmatrix}$ of \cref{we:linalg:cov} (grey;
an illustrative sample, seed 0, \texttt{tools/compute\_linalg.py}) and the four worked pixels
(teal). The ellipses are the points at whitened distance 1 and 2; the arrows are the eigenvectors
scaled by the square roots of their eigenvalues, 3.32 and 1. (b) The same points after
$\vect{z} = \mat{W}\vect{d}$: the ellipses have become circles, and the eigenvectors the
coordinate axes.}
\label{fig:linalg:whitening}
\end{figure}

\Cref{fig:linalg:whitening} shows the whole chapter in one picture. On the left, the cloud is a
tilted ellipse: long along $\vect{u}_1$, short along $\vect{u}_2$. On the right, after rotating
by $\mat{U}\T$ and dividing by $\sqrt{\lambda_k}$, it is round. Two points the same Euclidean
distance from the centre on the left can be very different distances on the right: a step of 1
unit along $\vect{u}_1$ becomes $1/3.32 = 0.30$ whitened units, while the same step along
$\vect{u}_2$ stays 1. Distance in the whitened picture is the \emph{Mahalanobis distance}, and
\cref{ch:rx} takes it from here.

\begin{keypoint}
Whitening is rotate onto the principal directions, then divide each by its standard deviation;
afterwards the background cloud is round and every direction counts equally.
\end{keypoint}

\subsection{Two practical details you will meet in the code}

\paragraph{Whitening with fewer directions.} If you keep only the top $d$ eigenvectors (the
first $d$ rows of $\mat{W}$), $\vect{z}$ has $d$ entries and is still white: its covariance is
the $d\times d$ identity, by the same proof with $\mat{U}$ replaced by its first $d$ columns
(\cref{ex:linalg:5}). The residual score of \cref{ch:residualmaha} and the PCA-RX baseline of
\cref{ch:localrx} do exactly this with $d = 30$.

\paragraph{A floor under the eigenvalues.} Dividing by $\sqrt{\lambda_k}$ is impossible when
$\lambda_k = 0$ and dangerous when it is tiny. Code that whitens therefore raises tiny eigenvalues
to a floor before the square root: \code{np.maximum(evals[order], 1e-12)} in
\code{app/utils/anomaly_detection/residual_mahalanobis.py} line 46, and
\code{np.maximum(noise_eigvals, 1e-10)} in \code{app/detectors/mnf_compression_detector.py}
line 321. Why tiny eigenvalues are dangerous, and not only zero ones, is the subject of
\cref{sec:rx:singular}.

\paragraph{Two whitening matrices.} The matrix $\mat{W} = \mat{\Lambda}^{-1/2}\mat{U}\T$ rotates
and then scales. Multiplying it on the left by $\mat{U}$ gives
$\mat{U}\mat{\Lambda}^{-1/2}\mat{U}\T$, which rotates, scales, and rotates back. It also whitens
(the extra $\mat{U}$ is a rotation, and $\mat{U}\mat{I}\mat{U}\T = \mat{I}$), and it is the
symmetric \emph{inverse square root} $\Cov^{-1/2}$, because multiplying it by itself gives
$\mat{U}\mat{\Lambda}^{-1}\mat{U}\T = \Cov\inv$. The MNF code builds this second form from the
noise covariance (\code{whiten = noise_eigvecs @ np.diag(1.0 / np.sqrt(noise_eigvals)) @
noise_eigvecs.T}, line 323); \cref{ch:localrx} explains why it does not matter which form you
use there.

\begin{confusions}
\confusion{Is the dot product the same as multiplying two vectors?}{No. Multiplying entry by
entry gives another vector. The dot product adds those products into a single number. The outer
product $\vect{d}\vect{d}\T$ is a third thing again: a whole matrix of products.}
\confusion{Why does projection divide by $\vect{u}^\top\vect{u}$?}{To make the answer independent
of the length of $\vect{u}$. Doubling $\vect{u}$ doubles $\vect{u}\T\vect{x}$ and quadruples
$\vect{u}\T\vect{u}$, so $t^\star$ halves and $t^\star\vect{u}$ is unchanged. With a unit vector
the division is by 1 and disappears.}
\confusion{Why $N-1$ in the covariance?}{Because the deviations are measured from the sample's own
mean and add to zero, so only $N-1$ of them are free. The code uses $N-1$ (\texttt{np.cov}). The
difference matters only for small samples, such as the local windows of \cref{ch:localrx} and the
tiny examples here.}
\confusion{Is an eigenvector unique?}{Only up to sign and length. If $\vect{u}$ is an
eigenvector, so is $-\vect{u}$ and $5\vect{u}$. We always take unit length; the sign is
arbitrary, and \texttt{np.linalg.eigh} may return either. Squared coordinates, and therefore every
score in this book, do not care about the sign.}
\confusion{Does PCA find the most interesting directions?}{It finds the directions of greatest
variance. Whether those are interesting depends on what produces the variance. If the noisiest
bands have the largest variance, PCA keeps the noise. \Cref{ch:localrx} builds MNF to fix exactly
that.}
\confusion{Whitening makes everything the same size, so does it destroy information?}{No. It is
an invertible change of coordinates (when all eigenvalues are kept and positive): you can undo it
with $\mat{U}\mat{\Lambda}^{1/2}$. What changes is which differences look large. That change is
the whole point.}
\end{confusions}

\begin{checkyourself}
\item\label{ex:linalg:1} Compute the angle in degrees between $(1, 2, 2)$ and $(2, 1, 2)$. Then
  project $(1, 2, 2)$ onto $(2, 1, 2)$ and check that the residual is orthogonal to $(2, 1, 2)$.
\item\label{ex:linalg:2} Three pixels in two bands: $(1, 2)$, $(3, 2)$, $(2, 5)$. Compute the
  mean, the deviations, the outer products and the covariance matrix (divide by $N - 1$). What is
  the correlation of the two bands?
\item\label{ex:linalg:3} Find the eigenvalues and unit eigenvectors of
  $\begin{pmatrix} 5 & 2\\ 2 & 2\end{pmatrix}$ by hand, and the explained variance ratio of the
  first component. Rebuild the matrix from $\sum_k \lambda_k\vect{u}_k\vect{u}_k\T$.
\item\label{ex:linalg:4} Using the matrix of \cref{ex:linalg:3}, write down the whitening matrix
  $\mat{W} = \mat{\Lambda}^{-1/2}\mat{U}\T$ and verify by multiplication that
  $\mat{W}\Cov\mat{W}\T = \mat{I}$.
\item\label{ex:linalg:5} Let $\mat{U}_d$ hold only the first $d$ eigenvectors of a $C\times C$
  covariance $\Cov$ as columns, and $\mat{\Lambda}_d$ the first $d$ eigenvalues. Show that
  $\vect{z} = \mat{\Lambda}_d^{-1/2}\mat{U}_d\T(\vect{x} - \muv)$ has covariance $\mat{I}_d$. Which
  fact about the columns of $\mat{U}$ does the proof use?
\item\label{ex:linalg:6} Show that $(\mat{A}\mat{B}\mat{C})\T = \mat{C}\T\mat{B}\T\mat{A}\T$ using
  \cref{eq:linalg:abt} twice, and explain in one sentence why $\mat{A}\Cov\mat{A}\T$ is symmetric
  whenever $\Cov$ is.
\end{checkyourself}

\nextchapter{ch:rx}
